% Experiment report. Fixed sections: see experiments/TEMPLATE/report.tex and experiments/README.md.
% Build: npm run experiment:report -- TEMPLATE
\documentclass{../polyexp}
\begin{document}
\input{generated}
\ExpHeader

\section{Question and hypothesis}
Stepping through replayed training games of experiment 24 (27 Nov 2026), Jakob marked seven failure modes:
``to the centre should score highest'' (all four first moves of a warrior scored the same, +3.61);
``one of them randomly moved to the edge'';
``a village was discovered by the last unit but stepping towards it now doesn't score higher'';
``this unit stands idle for some reason''; ``does nothing''; ``moves back to the corner''.
The replays show two causes. (1) Every feature describes what the unit can already see; under the fog all directions
look the same, ties are broken at random (edge moves) and from turn 1 every move has gain exactly 0, so the units pass.
(2) Travel distances count whole turns, so a single step toward a target often leaves the distance unchanged.

Jakob: ``I think we need to reward it for discovering the map\ldots'' The game already pays for this: the score gives
5 points per explored tile (\texttt{score\_of}, \texttt{core/src/state.rs}), but experiment 24's unit reward leaves the
score out.

\todo{jakob}{Confirm or change the prediction below.}
Behaviour count: the share of unit-turns that end idle although the unit still had options. In 12 replayed games of
experiment 24's linear run (iteration 800) it is 10\% in turns 1--10 and 34\% in turns 11--30.
Prediction (Claude's proposal): with $w = 0.5$ idle unit-turns in turns 11--30 fall from 34\% to under 20\%, villages
captured per game rise from 2.2 to at least 3, and the learner minus the champions'
mean improves by at least 300 points over the base (experiment 24 linear at iteration 300). $w = 1$ explores more but
captures fewer cities than $w = 0.5$. No improvement at any $w$ would count against the idea.

\section{Mechanism}
Every action of a learner unit earns
\[ r_{\text{explore}} = w \times (\text{tiles the learner's player has explored after the action} - \text{before}), \]
credited to the acting unit, on top of experiment 24's kill, capture and city-lost rewards. Because the greedy choice
adds each candidate's immediate reward to its gain, a move into the fog now beats a move back home. Knob
\texttt{--explore} (stars per tile, default 0 = experiment 24), \texttt{explore\_gain} and \texttt{credit} in
\texttt{engine/unitrl/src/main.rs}; the analysis shows it as the reward source ``explore''.

\section{Setup}
Three runs, $w \in \{0.25, 0.5, 1\}$ star per tile, sized against the other rewards: a village is worth its income
times the turns left (about 20--40 stars early), so at $w = 0.5$ revealing a typical move's 3--5 tiles is worth
1.5--2.5 stars, less than heading for a known village. In a 3-iteration smoke test at $w = 0.5$ warriors earned about
22 stars per game from exploring against about 40 from villages. Everything else as experiment 24's linear run
(feature set 1, $\gamma = 1$, Monte Carlo, $\varepsilon$ 0.3 $\to$ 0.05 over 200 iterations, 128 games per
iteration), but 300 iterations; the base is experiment 24's linear weights at iteration 300, measured on the same
200 evaluation seeds. Runs: \texttt{rl-unit-explore-025}, \texttt{rl-unit-explore-05}, \texttt{rl-unit-explore-1}.

\section{Results}
% WHEN: after the runs. The tables and plots below are generated from the measured data; add two to four sentences that read them.
\ExpResults

\todo{claude}{Read the tables: the best point, the size of the effect against the noise (about $\pm$1.2 points on a 6-model mean), the head-to-head.}

\section{Behaviour}
% WHEN: after the runs. WHO: both. One count the mechanism should change, plus what was seen in watched games.
\todo{claude}{The behaviour metric: one count the mechanism should change, with base and new value.}
\todo{jakob}{What you saw in the live games, if you watched any.}

\section{Interpretation}
% WHEN: after the runs. WHO: both. Why it worked or did not; was the prediction right?
\todo{claude}{Why the result came out this way, and whether the prediction held.}
\todo{jakob}{Your reading: agree, disagree, what is missing?}

\section{Threats to validity}
% WHO: Claude. What could make this result wrong or narrow: panel, map sizes, sample size, interactions with other knobs.
\todo{claude}{Two or three concrete threats.}

\section{Verdict and next step}
% WHO: Claude proposes by the keep rules (.claude/skills/model-experiment/SKILL.md); Jakob decides anything outside them.
\todo{claude}{Kept or not, under which keep rule, the champion afterwards, and the next experiment this suggests.}

\section{Notes for the paper}
% WHO: Jakob, with Claude's draft. What this experiment contributes: a claim, a figure, a negative result worth a line.
\todo{claude}{Draft: the one-sentence claim and which figure or table would carry it.}
\todo{jakob}{Is this worth a place in the paper? Where?}

\end{document}
